\chapter{异常}
\section{异常}
异常$\neq$错误

异常是在程序运行过程中出现的错误，而错误是程序的语法错误。

异常处理的一般流程是：检测到错误，引发异常，捕获异常，处理异常。

\begin{lstlisting}[style=demostyle,numbers=none]
# 第一种 try ... except
## 1.当try代码块中的代码报错时，
##  只有抛出与except关键字后异常类型相同的异常，才会执except代码块中的代码
## 2.如果except关键字后没有接异常类型，
##  那么任何异常都会触发except代码块中的代码逻辑
try:
    # 代码块
    pass
except Exception as e:
    # 异常处理
    pass

# 第二种 try ... except ... finally
try:
    # 代码块
    pass
except Exception as e:
    # 异常处理
    pass
finally:
    # 无论是否发生异常都会执行
    pass

# 第三种 try ... except ... else
try:
    # 代码块
    pass
except Exception as e:
    # 异常处理
    pass
else:
    # 无异常时执行
    pass

# 第四种 try ... except ... else ... finally
try:
    # 代码块
    pass
except Exception as e:
    # 异常处理
    pass
else:
    # 无异常时执行
    pass
finally:
    # 无论是否发生异常都会执行
\end{lstlisting}

\section{异常的种类}
\subsection{文件}
\begin{lstlisting}[style=pythonstyle,numbers=none]
try:
    f = open('test.txt', 'r')
    print(f.read())
except FileNotFoundError as e:
    print('文件不存在')
finally:
    f.close()
\end{lstlisting}

\subsection{除零}
\begin{lstlisting}[style=pythonstyle,numbers=none]
try:
    print(5 / 0)
except ZeroDivisionError as e:
    print('除零错误')
\end{lstlisting}

\subsection{输入异常}
\begin{lstlisting}[style=pythonstyle,numbers=none]
year = int(input('请输入年份：'))

try:
    if year % 4 == 0 and year % 100 != 0 or year % 400 == 0:
        print('闰年')
    else:
        print('平年')
except ValueError as e:
    print('输入错误')
\end{lstlisting}

\subsection{捕获多个异常}
\begin{lstlisting}[style=pythonstyle,numbers=none]
try:
    print(5 / 0)
    year = int(input('请输入年份：'))
except ZeroDivisionError as e:
    print('除零错误')
except ValueError as e:
    print('输入错误')
\end{lstlisting}

\begin{lstlisting}[style=pythonstyle,numbers=none]
try:
    print(5 / 0)
    year = int(input('请输入年份：'))
except (ZeroDivisionError, ValueError) as e: # 用一个元组来捕获多个异常
    print('输入错误')
\end{lstlisting}

\subsection{捕获所有异常}
\begin{lstlisting}[style=pythonstyle,numbers=none]
try:
    print(5 / 0)
    year = int(input('请输入年份：'))
except Exception as e:  # Exception是所有异常的基类
    print('发生异常')
\end{lstlisting}

\subsection{自定义异常返回信息——raise}
\begin{lstlisting}[style=pythonstyle,numbers=none]
try:
    raise NameError('自定义异常')
except NameError as e:
    print("My custom Exception")
\end{lstlisting}

\subsection{自定义异常}
\section{异常的处理}
通过\codeinline{try ... except}可以捕获异常，但如果捕捉到异常无法处理呢？此时最好的方式就是向上抛出异常。

异常处理的原则是处理可预见的异常，向上抛出当前无法处理的异常，最终的目的就
是方便程序编写者可以很轻松地定位出现异常的位置，以及出现异常的原因。

\begin{lstlisting}[style=demostyle,numbers=none]
# raise
raise [exceptionName('reason')]
## 中括号“[]”括起的部分是可选参数，
## exceptionName的作用就是抛出指定类型的异常，
## reason则是抛出异常时附带的说明信息。
\end{lstlisting}

\begin{lstlisting}[style=pythonstyle,numbers=none]
def error_fun():
    try:
        name = input('Enter your name: ')
        if name.isspace():
            raise ValueError('Name cannot be empty')
        if name.isdigit():
            raise ValueError('Name cannot be a number')
        
        name = name + 123

    # 无法处理本层抛出的ValueError,会抛出该错误到上一层
    except ValueError as e:
        print(f'错误类型: {e}') # f的作用是将{}中的内容替换为变量的值
    else:
        print('Name:', name)

# error_fun()层的上一层是__main__，所以抛出的错误会被__main__处理
try:
    error_fun()
except ValueError as e:
    print(f'主函数抛出异常: {e}, Please try again.')
    error_fun()
\end{lstlisting}


